\chapter{Theoretical Background}\label{ch:theoretical}

This chapter provides the theoretical foundation for
Chapter~\ref{ch:proposed}: an axis-based view of linear-time
recurrent expressivity, a unified recurrence template for the
three matrix backbones, the LION bidirectional adaptation with
its decay-class taxonomy, the three primitives the proposed
mechanisms build on, and the formal complexity-class bounds that
delimit the diagonal class within which the thesis operates.

\section{Axis decomposition of linear-time RNN expressivity}\label{sec:axes}

In this thesis, \textit{expressivity} refers to the formal
sequence-transformation capacity of the recurrent state-update
operator, in the tradition of the formal-expressivity results
of~\citet{sarrof2024expressive} and the architectural-capacity
framing of~\citet{yang2024deltanet} and~\citet{siems2025deltaproduct}
on DeltaNet and DeltaProduct, rather than aggregate benchmark
performance.

The expressivity of a sequence-modelling architecture is sometimes
treated as a single-scalar property, with one architecture said to
be \textit{more expressive} than another. This thesis adopts a
different position. The expressivity of causal linear-time
recurrent architectures does not collapse to a single scalar; it
decomposes into conceptually distinct but empirically interacting
axes, each exercised by a different task structure, and a
mechanism that extends the function class along one axis does not
generally convert into measured gain on a task whose structural
prior does not exercise that axis. The decomposition is the
conceptual scaffold around which the rest of the thesis is
organised.

Three axes are central to the thesis. Axes one and two are
directly exercised by the matrix's primary probes; axis three
(non-Abelian state-tracking) is out of scope for the
diagonal-class mechanisms proposed in this thesis but is
bracketed by the formal complexity-class bound of
Section~\ref{sec:bounds}. Two additional axes are mentioned for
completeness and for positioning the discovery-phase archive that
preceded the matrix, but are not exercised on the matrix's
primary probes.

The first axis concerns the architecture's capacity to encode
short-range temporal structure. The relevant patterns include
phoneme-to-syllable transitions in speech, formant trajectories
in voiced segments, $n$-gram structure in language, and
adjacent-pixel structure in vision. The deficit is the absence
of explicit local mixing before the recurrent state update;
Section~\ref{sec:primitives:msdc} gives the technical
formulation. The diagnostic property of an axis-one task is that
improvements in short-range mixing translate into reductions of
the primary metric, which on the thesis's ASR probes is the
character error rate at post-subsampling sequence lengths
$T \le 500$. Linear-time architectures exhibit characteristic
failure modes on axis-one tasks, most visibly through the
scaling difficulties of Vision Mamba~\cite{zhu2024visionmamba}
on image classification beyond modest parameter counts and
through the documented gap of Audio
Mamba~\cite{erol2024audiomamba} on speech recognition in the
absence of explicit local mixing.

The second axis concerns the architecture's capacity to store and
retrieve key-value associations when many of the keys are similar.
Recurrent linear-time architectures compress the history into a
bounded state of dimension $d$: when the number of stored pairs $N$
approaches or exceeds $d$, similar keys overlap in the state and
become difficult to disentangle at readout. Softmax attention, by
contrast, performs a non-collapsing pairwise comparison and does not
exhibit the same interference pattern at the same scale. The
diagnostic property of an axis-two task is the structured saturation
of recall accuracy as $N/d$ grows. The Multi-Query Associative
Recall benchmark of~\citet{arora2023zoology}, designed exactly for
this property, sweeps the sequence length $T \in \{64, 256, 1024\}$
with the number of stored pairs growing in proportion, and is the
cleanest available probe for axis two in the linear-time recurrent
literature.

The third axis concerns the architecture's capacity to track
discrete state across arbitrary sequence length, including the word
problems of finite groups. Diagonal linear-time recurrences,
encompassing all three backbones of Section~\ref{sec:backbones} and
the block-complex extension of Section~\ref{sec:primitives:dho}, are
bounded above by the constant-depth threshold-circuit class
$\mathsf{TC}^{0}$~\cite{merrill2024illusion}. The result
of~\citet{sarrof2024expressive} sharpens this further: complex-
diagonal state-space models of depth $L$ cannot represent the word
problem for any non-Abelian group whose subnormal Abelian-factor
series exceeds length $L$, regardless of width. Axis three is
therefore bracketed by formal complexity theory and is out of scope
for the mechanisms proposed in this thesis. The diagonal-plus-low-rank lift of
DeltaNet~\cite{schlag2021deltanet, yang2024deltanet},
DeltaProduct~\cite{siems2025deltaproduct}, and
RWKV-7~\cite{peng2025rwkv7} is the principal route by which a
recurrent linear-time architecture crosses the formal class
boundary, and is referenced here to mark the boundary rather than
exercised. Section~\ref{sec:bounds} returns
to this axis in detail.

Two further axes appear in the related-work map of the thesis but
are not central to the proposed mechanisms. The fourth axis,
long-range information flow at sequence lengths far beyond the
chunk-local window of standard linear-time recurrent architectures,
has its own line of work that includes log-linear
attention~\cite{guo2025loglinear} and chunk-level sparse retrieval
mechanisms~\cite{hou2025rwkvx, hu2025hsa}; these address regimes
that exceed the operating range of the discovery probe used in this
thesis ($T \le 500$). The fifth axis, content-adaptive feature-map
richness through polynomial value lifts and channel-side
cross-products, has been explored across a saturated cluster of
polynomial-mixer parametrisations~\cite{mongaras2025expressiveness, picard2026pom, hammoud2025dont, zhang2025hla}.
Mechanisms targeting axes four and five were explored in the
discovery phase of the thesis on backbones whose primary task
exercises axes one or two. The empirical matrix of
Chapter~\ref{ch:experiments} populates the positive-transfer half
of the criterion only and does not report a converse-direction
catalogue; the discovery-phase archive remains the source for the
engaged-but-null evidence that the criterion's full statement
predicts.

Across the three primary axes, the thesis observes a single
empirical regularity that the synthesis chapter formalises as a
predictive criterion: larger residual deficits along the axis a
mechanism targets tend to correspond to larger absolute gains,
and no measured gain is expected when the mechanism extends the
function class along an axis the task does not exercise.
Chapter~\ref{ch:proposed} instantiates this criterion in three
mechanisms, and Chapter~\ref{ch:experiments} tests it across
backbones and probes.

\begin{table}[t]
\centering
\small
\caption[Axes of expressivity for causal linear-time RNNs]{Axes
of expressivity considered in this thesis, with the diagnostic
property that distinguishes each axis and the task family that
exercises it. Axes one and two are exercised by the matrix's
primary probes; axis three is bracketed by the formal
complexity-class bound of Section~\ref{sec:bounds}; axes four
and five are mentioned for completeness.}
\label{tab:axes}
\setlength{\tabcolsep}{4pt}
\renewcommand{\arraystretch}{1.15}
\begin{tabularx}{\linewidth}{@{}p{0.4cm} p{3.0cm} X p{3.0cm}@{}}
\toprule
\textbf{\#} & \textbf{Axis} & \textbf{Diagnostic property} & \textbf{Task / probe} \\
\midrule
1 &
Short-range temporal hierarchy &
Improvements in explicit short-range mixing translate into
reductions of the primary metric &
ASR (LibriSpeech, Common Voice) \\
\addlinespace
2 &
Associative memory under interference &
Recall accuracy saturates as the number of stored key-value pairs
$N$ approaches the state dimension $d$ &
MQAR with $T \in \{64, 256, 1024\}$ \\
\addlinespace
3 &
State-tracking, non-Abelian permutation &
Architecture position relative to the $\mathsf{NC}^{1}$ /
$\mathsf{PNC}^{1}$ class boundary &
$S_{n}$ permutation, Dyck-$k$ ($k \ge 2$); out of scope \\
\addlinespace
4 &
Long-range information flow &
Capacity at $T$ far beyond the chunk-local window &
NIAH, BABILong; out of scope at $T \le 500$ \\
\addlinespace
5 &
Content-adaptive feature-map richness &
Channel-side polynomial-lift expressivity &
Polynomial-mixer benchmarks; discovery-phase only \\
\bottomrule
\end{tabularx}
\end{table}

\section{Three causal backbones}\label{sec:backbones}

Causal RWKV-6, Mamba-2, and Linear Attention are presented in the
literature with distinct notations and differing emphases on
continuous-time, recurrent, or kernel-attention interpretations. This
section formalises the three within a single unified discrete-time
recurrence template, which makes their structural differences visible
at a glance and allows the rest of the thesis to refer to the three
as variants of a single object. Each backbone's canonical form is
reproduced as it appears in the source paper that introduced the
architecture, with notation alignment performed only to make the
unified template readable.

The most general discrete-time linear recurrence considered in this
thesis takes the form
\begin{equation}\label{eq:unified}
h_{t} = A_{t} \odot h_{t-1} + B_{t} \, x_{t}, \qquad y_{t} = C_{t}^{\top} h_{t},
\end{equation}
where $h_{t}$ is the hidden state, $A_{t}$ is an elementwise (or
diagonal-structured) transition operator at step $t$, $B_{t}$ is the
input-write operator, $C_{t}$ is the readout direction, and $\odot$
denotes the elementwise product. The notation matches the formal SSM
definition of~\citet{sarrof2024expressive} and the discrete S6 form
of~\citet{dao2024transformers}, in both of which the bar notation
distinguishing continuous and discrete parameters is intentionally
dropped to focus on the discrete dynamics. The three backbones used
in the experimental matrix instantiate
equation~(\ref{eq:unified}) with different choices of $A_{t}$,
$B_{t}$, and the structure of $h_{t}$. The structural axis on which
they differ is primarily the transition operator $A_{t}$, which
varies from the trivial accumulator $A_{t} = I$ of Linear Attention,
through the per-channel data-dependent diagonal of RWKV-6, to the
input-modulated scalar-identity form of Mamba-2 SSD.

\subsection{Linear Attention as a pure accumulator}\label{sec:backbones:la}

The Linear Attention architecture of~\citet{katharopoulos2020transformers}
expresses causal attention as a recurrence with bounded state. With
$\phi(\cdot)$ a positive feature map applied to the keys and queries,
the recurrence is
\begin{equation}\label{eq:la-state}
S_{t} = S_{t-1} + \phi(K_{t}) V_{t}^{\top}, \qquad
z_{t} = z_{t-1} + \phi(K_{t}),
\end{equation}
\begin{equation}\label{eq:la-output}
y_{t} = \frac{\phi(Q_{t})^{\top} S_{t}}{\phi(Q_{t})^{\top} z_{t}},
\end{equation}
where $S_{t} \in \mathbb{R}^{d_{K} \times d_{V}}$ is the matrix-valued
state, $z_{t} \in \mathbb{R}^{d_{K}}$ is the normalisation
accumulator, and $K_{t}, V_{t}, Q_{t}$ are the per-token key, value,
and query projections of the input $x_{t}$. The state is a pure
accumulator: each token write is added to the running sum without any
attenuation of earlier writes.

Within the unified template of equation~(\ref{eq:unified}),
Linear Attention corresponds to the trivial transition
$A_{t} = I$, an input-write operator $B_{t} = \phi(K_{t})$ that
produces a rank-one outer-product update with $V_{t}$, and a
separable readout direction $C_{t} = \phi(Q_{t})$ applied with a
normalisation step over $z_{t}$.

\subsection{RWKV-6 with per-channel data-dependent decay}\label{sec:backbones:rwkv6}

The RWKV-6 architecture, introduced as Finch
in~\citet{peng2024eaglefinch}, builds on the matrix-state RWKV-5/6
family with a per-channel data-dependent decay. The Time Mix block
computes a state update and a readout in two coupled equations:
\begin{equation}\label{eq:rwkv6-state}
s_{t} = \mathrm{diag}(w_{t}) \cdot s_{t-1} + k_{t}^{\top} v_{t},
\end{equation}
\begin{equation}\label{eq:rwkv6-readout}
\mathrm{wkv}_{t} = s_{t-1} + \mathrm{diag}(u) \cdot k_{t}^{\top} v_{t},
\end{equation}
where $s_{t} \in \mathbb{R}^{(D/h) \times (D/h)}$ is the per-head
matrix-valued state, $w_{t} \in \mathbb{R}^{D/h}$ is the per-channel
data-dependent decay vector, $u \in \mathbb{R}^{D/h}$ is a static
bonus vector applied to the current rank-one outer product, and
$k_{t}^{\top} v_{t}$ is the rank-one write at step $t$. Note the
recurrence order: the readout $\mathrm{wkv}_{t}$ uses the previous
state $s_{t-1}$ together with the bonus on the current write, and the
state then advances to $s_{t}$ for the next step.

The decay vector $w_{t}$ is data-dependent. It is produced by a
data-dependent linear interpolation of $x_{t}$ and $x_{t-1}$,
followed by a low-rank LoRA-style projection and a stabilising
double-exponential map:
\begin{equation}\label{eq:rwkv6-decay}
d_{t} = \mathrm{lora}_{d}\!\big(\mathrm{ddlerp}_{d}(x_{t}, x_{t-1})\big),
\qquad w_{t} = \exp\!\big({-}\exp(d_{t})\big),
\end{equation}
where $\mathrm{ddlerp}_{\ast}(a, b) = a + (b - a) \odot \mathrm{lora}_{\ast}\!\big(a + (b - a) \odot \mu_{x}\big)$
is the Finch data-dependent linear-interpolation operator and
$\mathrm{lora}_{\ast}(x) = \lambda_{\ast} + \tanh(x A_{\ast}) B_{\ast}$
is its low-rank residual core, with $\mu_{x}$ and $\lambda_{\ast}$
trainable vectors and $A_{\ast}, B_{\ast}$ trainable low-rank
projections. The double-exponential map ensures $w_{t} \in (0, 1)^{D/h}$
without requiring an explicit clamp.

Within the unified template of equation~(\ref{eq:unified}),
RWKV-6 corresponds to a transition
$A_{t} = \mathrm{diag}(w_{t})$ with $w_{t}$ per-channel
data-dependent through equation~(\ref{eq:rwkv6-decay}), an
input-write that produces the rank-one outer product
$k_{t}^{\top} v_{t}$, and a readout that includes the static
bonus $\mathrm{diag}(u)$ on the current write in addition to the
previous-state component $s_{t-1}$.

\subsection{Mamba-2 with input-modulated scalar transitions}\label{sec:backbones:mamba2}

Mamba-2, introduced in~\citet{dao2024transformers}, refines the
selective state-space model of~\citet{gu2023mamba} with the
structured state-space duality (SSD) framework. The general
selective recurrence has the form
\begin{equation}\label{eq:mamba2-s6}
h_{t} = A_{t} h_{t-1} + B_{t} x_{t}, \qquad y_{t} = C_{t}^{\top} h_{t},
\end{equation}
where $A_{t}, B_{t}, C_{t}$ are the discrete (per-step) parameters of
the SSM. The SSD structural restriction sets $A_{t} = a_{t} I$ for a
scalar $a_{t} \in \mathbb{R}$ and the identity matrix $I$, which is
the formal step that connects Mamba-2 to masked kernel attention with
one-semiseparable masks~\cite[Section~3]{dao2024transformers}.
Substituting the restriction into equation~(\ref{eq:mamba2-s6}) yields
\begin{equation}\label{eq:mamba2-ssd}
h_{t} = a_{t} \cdot h_{t-1} + B_{t} x_{t}, \qquad y_{t} = C_{t}^{\top} h_{t}.
\end{equation}

The scalar $a_{t}$ is parameterised by an input-dependent step size
$\Delta_{t}$ and an underlying continuous parameter $\mathring{A}$
through a zero-order-hold discretisation
of~\citet[eq.~zoh]{gu2023mamba}: $a_{t} = \exp(\Delta_{t} \mathring{A})$.
The selective step size $\Delta_{t}$ is itself input-dependent through
a low-rank linear projection followed by a softplus activation:
\begin{equation}\label{eq:mamba2-delta}
\Delta_{t} = \mathrm{softplus}\!\Big(\mathsf{P} + \mathrm{Broadcast}_{D}\!\big(\mathrm{Linear}_{1}(x_{t})\big)\Big),
\end{equation}
where $\mathsf{P}$ is a learnable per-channel bias parameter and
$\mathrm{Linear}_{1}(x_{t})$ projects the input to a scalar that is
broadcast across the channel dimension before being added to the bias.
The discrete input-write parameter $B_{t}$ and readout direction
$C_{t}$ are similarly produced by input-dependent linear projections
$s_{B}(x_{t}) = \mathrm{Linear}_{N}(x_{t})$ and $s_{C}(x_{t}) = \mathrm{Linear}_{N}(x_{t})$~\cite{gu2023mamba}.

Within the unified template of equation~(\ref{eq:unified}),
Mamba-2 corresponds to a transition $A_{t} = a_{t} I$ with $a_{t}$
per-token data-dependent through the selective step size
$\Delta_{t}$ of equation~(\ref{eq:mamba2-delta}), and
input-dependent $B_{t}$, $C_{t}$ through low-rank projections of
$x_{t}$.

\subsection{Architecture-deficit map}\label{sec:backbones:deficit}

The three instantiations differ along the structural axis of the
transition operator and along the input-write geometry, and each
leaves a different residual deficit uncovered. Linear Attention has
no native decay and no explicit local mixing, and the unbounded
accumulator implies that interference between similar keys
accumulates without any architectural correction. RWKV-6 introduces
per-channel data-dependent decay, which addresses the absent-state-
attenuation deficit on Linear Attention's terms, but its rank-one
outer-product writes provide no structural mechanism for active
interference cancellation between correlated keys. Mamba-2 introduces
an input-dependent scalar transition $a_{t}$ that selects how
strongly each token write persists, partially closing the
content-adaptive routing deficit, but the SSD scalar-identity
restriction means the transition does not differentiate per-channel
decay rates, and the architecture has no native interference-
cancellation mechanism either. None of the three backbones has an
explicit input-side local-mixing operator at the post-subsampling
scale relevant to short-range temporal hierarchy.

The proposed mechanisms of Chapter~\ref{ch:proposed} target these
residual deficits one at a time. The Multi-Scale Depthwise
Convolution mechanism, which builds on the input-side local-mixing
primitive of Section~\ref{sec:primitives:msdc}, addresses the
short-range temporal hierarchy deficit shared by all three backbones.
The Damped Harmonic Oscillator mechanism, which builds on the
block-complex transition primitive of
Section~\ref{sec:primitives:dho}, targets the damped-oscillation
sub-axis on architectures whose native transition is real-diagonal
or trivial. The Chunked Value Decorrelation mechanism, which builds
on the key-similarity preconditioning primitive of
Section~\ref{sec:primitives:cvd}, addresses the absent
interference-cancellation deficit on the linear-time recurrent
substrate.

\section{Bidirectional adaptation via LION}\label{sec:lion}

The three causal linear-time backbones formalised in
Section~\ref{sec:backbones} propagate state in one direction, which
suits streaming and autoregressive workloads but discards future
context that is available in offline settings such as full-utterance
speech recognition. The related-work chapter surveyed three
structurally distinct strategies for adapting an existing causal
architecture to bidirectional input access: a mechanism-level
parallel attention, an operator-level forward-plus-backward scan, and
an ordering-level per-layer scan alternation. This thesis adopts the
mechanism-level strategy and instantiates the LION framework
of~\citet{afzal2025linear} as a unified parallel adaptation across
the three backbones. The choice is motivated by the structural
property that the parallel form inherits the causal mechanism
vocabulary directly and does not double the serial compute per layer;
an empirical comparison against the operator-level alternatives of
Vision-RWKV's Bi-WKV scan~\cite{duan2024visionrwkv} and Vision
Mamba's forward-and-backward block pairing~\cite{zhu2024visionmamba}
is reserved for the experimental chapter.

The mathematical construction of LION starts from the observation
that causal linear attention, once written as a matrix-vector
operation rather than a per-token recurrence, takes the form
$\mathbf{Y} = \mathrm{Scale}\!\big((\mathbf{Q} \mathbf{K}^{\top}) \odot \mathbf{M}^{C}\big) \mathbf{V}$,
where $\mathbf{M}^{C}$ is a lower-triangular causal mask whose
off-diagonal entries encode the cumulative decay between any pair of
positions. Writing the entry of $\mathbf{M}^{C}$ in row $i$ and column
$j$ as $\mathbf{M}^{C}_{ij} = \prod_{k = j+1}^{i} \lambda_{k}$ for
$i > j$ makes explicit that the causal mask is a product of per-step
decay factors $\lambda_{k}$. The bidirectional generalisation drops
the lower-triangular constraint and uses a symmetric mask whose
entry between positions $i$ and $j$ equals the product of decay
factors between them in either direction:
\begin{equation}\label{eq:lion-mask}
\mathbf{M}_{ij} = \begin{cases}
\prod_{k = j}^{i - 1} \lambda_{k} & i > j, \\
1 & i = j, \\
\prod_{k = i + 1}^{j} \lambda_{k} & i < j.
\end{cases}
\end{equation}
The Full Linear Attention output is then
\begin{equation}\label{eq:lion-fullatt}
\mathbf{Y} = \mathrm{Scale}\!\big((\mathbf{Q} \mathbf{K}^{\top}) \odot \mathbf{M}\big) \mathbf{V},
\end{equation}
where $\odot$ denotes the elementwise product and $\mathrm{Scale}(\cdot)$
is a row-wise normalisation step. The mask $\mathbf{M}$ admits a
rank-one semi-separable factorisation~\cite{afzal2025linear} that
allows efficient construction via cumulative sums in log-space, which
keeps the parallel implementation tractable at large sequence lengths
while paying the quadratic memory footprint of the bidirectional
access pattern.

The structure of $\lambda_{k}$ defines a taxonomy of LION variants.
Three named instantiations appear in the LION paper, each
corresponding to a distinct causal backbone via the natural
parallel-form mapping:
\begin{itemize}
    \item \textbf{LION-LIT}: $\lambda_{k} = 1$ for all $k$. The mask
      becomes the all-ones matrix and
      equation~(\ref{eq:lion-fullatt}) reduces to the bidirectional
      analogue of Vanilla Linear
      Attention~\cite{katharopoulos2020transformers}. No decay.
    \item \textbf{LION-S}: $\lambda_{k} = \sigma(W x_{k} + b)$, where
      $\sigma$ is the sigmoid activation (with a shifted SiLU
      pre-activation in the reference implementation), $W$ and $b$
      are learnable, and the decay is per-token data-dependent.
      The variant is the bidirectional extension of Gated Random
      Feature Attention and is described as inspired by the
      input-dependent selectivity of Mamba-2.
    \item \textbf{Learnable fixed decay}: $\lambda_{k} = \lambda$
      for a single learnable scalar shared across positions, which
      makes $\mathbf{M}$ a Kac-Murdock-Szego matrix. This is the
      natural parallel form for architectures with a single global
      decay rate.
\end{itemize}

The per-backbone parallel-form mapping (RWKV-6 to LION-S,
Mamba-2 to LION-S, Linear Attention to LION-LIT, with a
controlled LION-S deployment on Linear Attention) is detailed
in Section~\ref{sec:proposed:lion}. The structural reason why
both LION-LIT and LION-S are reported on Linear Attention is
examined in the remainder of this section.

The role of decay in the bidirectional setting is more constraining
than in the causal setting. In the causal recurrence, the unbounded
accumulator of Linear Attention is acceptable because the readout
normalisation $z_{t} = \sum_{j \le t} \phi(K_{j})$ grows in proportion
to the state, and the ratio in equation~(\ref{eq:la-output}) remains
bounded. In the parallel bidirectional setting, the row sums of the
unmasked $\phi(\mathbf{Q}) \phi(\mathbf{K})^{\top}$ block grow as
$O(T \cdot d_{K})$, and the absence of a decay structure on
$\mathbf{M}$ means the value tensor enters the output proportional to
a sum that scales linearly in sequence length without architectural
attenuation. The LION-LIT instantiation remains a correct definition,
but a downstream mechanism that operates on this output, such
as the value-decorrelation primitive of
Section~\ref{sec:primitives:cvd}, loses its conditioning
guarantee when applied to a substrate whose accumulation is
unbounded. The matrix therefore reports both LION-LIT (the
natural parallel-form mapping; controlled baseline) and LION-S
(per-token sigmoid decay added) on Linear Attention; the
empirical outcome is reported in Chapter~\ref{ch:experiments}.

\section{Mathematical primitives for mechanism design}\label{sec:primitives}

Chapter~\ref{ch:proposed} introduces three proposed mechanisms designed
to address structural deficits of the three causal linear-time backbones
formalised in Section~\ref{sec:backbones}. Each proposed mechanism
builds on a mathematical primitive that has been previously introduced
in the literature, in a setting that differs from the one this thesis
exercises. This section presents the three primitives in their original
formulations and defers the linear-time-recurrent instantiations to the
proposed-solution chapter.

The three primitives are: the input-side multi-scale local-mixing
pattern, independently rediscovered across vision, audio, and language
settings; the block-complex transition with Rayleigh dissipation,
drawn from analytical mechanics and Lie-group theory; and the
key-similarity preconditioning operator introduced for softmax
attention by Duvvuri et al.~\cite{duvvuri2026lucid}. They underpin, respectively, the
Multi-Scale Depthwise Convolution, Damped Harmonic Oscillator, and
Chunked Value Decorrelation mechanisms of Chapter~\ref{ch:proposed}.

\subsection{Multi-scale depthwise dilated convolution}\label{sec:primitives:msdc}

The structural deficit is shared across the linear-time recurrent
family: each token sees only its own embedding plus the running
recurrent state, and that state cannot encode short-range structure as
efficiently as an explicit local operator can. The deficit has been
documented separately in vision through the scaling difficulties of
Vision Mamba~\cite{zhu2024visionmamba} beyond modest parameter counts,
in audio classification, and in long-context language modelling. The
literature converges on a single functional pattern: an input-side
local operator that gives every token explicit access to a short-range
neighbourhood before the recurrent core processes the sequence.

Four instances of the pattern appear across the linear-time literature,
differing in syntactic form rather than in function.
Vision-RWKV~\cite{duan2024visionrwkv} introduces a quad-directional
shift, named Q-Shift, that splits the channel dimension into four
groups and shifts each group by one position in one of $\{$up, down,
left, right$\}$ before concatenation:
\begin{equation}\label{eq:qshift-shift}
\begin{split}
X^{\dagger}[h, w] = \mathrm{Concat}\big(
&X[h{-}1,\, w,\, 0{:}C/4],\; X[h{+}1,\, w,\, C/4{:}C/2], \\
&X[h,\, w{-}1,\, C/2{:}3C/4],\; X[h,\, w{+}1,\, 3C/4{:}C]\big)
\end{split}
\end{equation}
\begin{equation}\label{eq:qshift-mix}
\mathrm{Q\text{-}Shift}(X) = X + (1 - \mu) \, X^{\dagger},
\end{equation}
with $\mu$ a learnable per-channel scalar that controls the residual
mixing strength. The connectivity pattern itself is fixed; only the
mixing weight is learned.

LiT~\cite{wang2025lit} attaches a depthwise convolution to the value
path of a linear attention layer,
\begin{equation}\label{eq:lit-vconv}
V \leftarrow \mathrm{DWConv}_{k}(V),
\end{equation}
introducing a single learnable kernel of size $k$ on each channel.
Compared with Q-Shift, the operator is fully learned but operates at a
single dilation.

The non-attention language model of Kiruluta et
al.~\cite{kiruluta2025breaking} generalises the fixed-shift and
single-convolution forms to a parallel bank of dilated convolutions:
\begin{equation}\label{eq:multi-resolution}
Z_{k} = \mathrm{Conv1d}(x,\, \mathrm{dilation} = d_{k}), \qquad k \in \{1, \ldots, K\},
\end{equation}
with the parallel branches combined typically through summation or
concatenation. The dilation set is described in the source as powers
of two, such as $1, 2, 4$, and so on, and the specific combination
scheme is left unspecified beyond the requirement that the parallel
branches feed back into the residual stream.

AudioRWKV~\cite{xiong2025audiorwkv} attaches a two-dimensional
depthwise-separable ConvShift, instantiated as a learnable variant of
Q-Shift, to the RWKV-7 backbone. The operator builds a local residual
that feeds into the decay, key, and value paths of the time-mix
block, and the paper's ablation reports a positive contribution from
the ConvShift in addition to the contribution from the underlying
Q-Shift.

The four instances differ on three orthogonal axes. The first is
whether the local connectivity pattern is fixed (Q-Shift, ConvShift)
or learned (LiT, multi-resolution). The second is whether multiple
spatial scales are aggregated in parallel (multi-resolution,
ConvShift) or a single scale is used (Q-Shift, LiT). The third is the
path on which the operator is applied: the input side, the value
path, or all three of the key, value, and decay paths. The common
functional pattern across the four is the introduction of an explicit
short-range mixing operator on top of the linear-time recurrent core,
restoring local inductive bias that the recurrence alone cannot
express. Chapter~\ref{ch:proposed} introduces the specific
instantiation used in this thesis, including the dilation schedule,
kernel size, mixing-weight parameterisation, and an initialisation
correction that arose during the discovery phase of the experimental
matrix.

\subsection{Block-complex transitions and Rayleigh dissipation}\label{sec:primitives:dho}

The transition operator of a linear-time recurrent architecture
determines the function class the hidden state can express across
time. RWKV-6's per-channel decay and Linear Attention's pure
accumulator $A_{t} = I$ both sit inside the real-diagonal class:
state channels evolve independently with no representation of
phase, and the function class cannot represent oscillatory
dynamics. This restricts the architecture's capacity to model
phenomena such as formant trajectories in voiced speech,
rotational trajectories in handwriting and motion, and other
processes whose state lives on a circle rather than a half-line.

The smallest natural extension of the real-diagonal class that
admits damped oscillation while remaining inside the
diagonalisable matrix family is the block-complex transition.
The standard reference for the rotation group $SO(2)$ and its
embedding into matrix Lie-group theory
is~\citet[Chapter 1]{hall2015lie}. Replacing each pair of state
channels with a two-vector that is rotated by an angle $\theta$
and scaled by a damping factor $e^{-\lambda}$ yields the
per-block transition
\begin{equation}\label{eq:block-complex}
G_{b} = e^{-\lambda} \, R(\theta), \qquad R(\theta) = \begin{pmatrix} \cos\theta & -\sin\theta \\ \sin\theta & \cos\theta \end{pmatrix}.
\end{equation}
Stacking $K/2$ such two-by-two blocks produces a block-diagonal
transition in $SO(2)^{K/2} \times \mathbb{R}_{+}^{K/2}$. Each
block has eigenvalues $r e^{\pm i \theta}$ with $r = e^{-\lambda}$,
placing the transition inside the complex-diagonal SSM class
characterised formally by~\citet{sarrof2024expressive}. The
block-complex form is therefore a structured subset of the
complex-diagonal class and inherits the formal-language
properties of that class as recalled in
Section~\ref{sec:bounds}.

Allowing the rotation angle $\theta$ to grow without further
structural constraint produces phase drift across many recurrent
steps and a sensitivity of the late-sequence state to the exact angle
distribution at early positions. The standard physical prior that
regulates this regime is the Rayleigh dissipation function from
analytical mechanics~\cite[Sections 1.5 and 2.5]{goldstein2002classical},
which couples damping to the squared angular velocity of an
oscillator. The discrete-step analogue, applied to the per-block
damping, takes the form
\begin{equation}\label{eq:rayleigh}
\lambda_{\mathrm{eff}} = \lambda_{\mathrm{raw}} + \eta \cdot \theta^{2},
\end{equation}
with $\eta \ge 0$ a viscosity coefficient. Higher-frequency rotations
damp more rapidly, and the rotation budget becomes self-regulating: a
clip on $|\theta|$ functions as a soft constraint rather than a hard
wall, and the long-sequence behaviour of the recurrent state remains
controlled. Setting $\eta = 0$ recovers the undamped block-complex
transition exactly, which makes the construction zero-initialisable:
the Rayleigh coupling reduces to the uncoupled block-complex form at
step zero of training and engages only as $\eta$ moves under the
gradient.

Chapter~\ref{ch:proposed} introduces the data-dependence
parameterisation of $\theta$ and the depth-graded rotation-budget
schedule used in this thesis as the Damped Harmonic Oscillator
mechanism, together with its three architecture-specific deployments.

\subsection{Key-similarity preconditioning and value decorrelation}\label{sec:primitives:cvd}

Recurrent linear-time architectures compress the history into a
bounded state. When several stored keys are similar, their associated
values overlap in the state and become difficult to disentangle at
readout. The same interference pattern arises in softmax attention at
long context lengths, where correlated keys cause the softmax
distribution to diffuse probability mass across multiple ostensibly
relevant tokens, degrading retrieval quality. The two settings differ
in substrate, but share the underlying problem: a representation
constructed from key-value pairs accumulates interference proportional
to the correlation among the keys.

A direct algebraic correction to this interference appears in the
LUCID Attention mechanism of~\citet{duvvuri2026lucid}, which targets
long-context retrieval inside a quadratic-complexity Transformer.
LUCID constructs a preconditioner from key-similarity statistics in
the reproducing kernel Hilbert space induced by the exponential
kernel:
\begin{equation}\label{eq:lucid-base}
P = M \circ \exp(K K^{\top}),
\end{equation}
where $M$ is the causal mask, $\circ$ denotes the elementwise
(Hadamard) product, and the exponential is applied entrywise. The
preconditioned values $P^{-1} V$ are then combined with standard
softmax attention weights to produce the LUCID attention output:
\begin{equation}\label{eq:lucid-output}
O = \mathrm{softmax}(Q K^{\top} / \sqrt{d}) \cdot P^{-1} V.
\end{equation}
For numerical stability and proper variance scaling, the
implementation reported in the LUCID paper applies a row-wise RMS
normalisation to the keys, denoted $K_{\mathrm{RN}}$, and a
$-\sqrt{d}$ shift inside the exponential. The resulting preconditioner
is unit-diagonal by construction:
\begin{equation}\label{eq:lucid-rn}
P = M \circ \exp\!\left( \frac{K_{\mathrm{RN}} K_{\mathrm{RN}}^{\top}}{\sqrt{d}} - \sqrt{d} \right),
\end{equation}
which improves the conditioning of $P$ at long context lengths and
admits a direct triangular solve for $P^{-1} V$.

The construction has a clean interpretation as a closed-form
delta-rule update in the exponential-kernel reproducing kernel
Hilbert space: applying $P^{-1}$ to $V$ removes from each value the
components that project onto the directions of earlier keys under the
key-similarity geometry. This is structurally analogous to the
rank-one Householder erasure introduced by
DeltaNet~\cite{schlag2021deltanet,yang2024deltanet}, where each
token's update subtracts an explicit rank-one component along the
current key, but realised in closed form across the entire sequence
rather than online per token.

The substrate of LUCID is softmax attention. The mechanism preserves
the $O(N^{2} d)$ complexity of standard attention, its target task
domain is multi-needle long-context retrieval, and the
$\mathrm{softmax}(Q K^{\top} / \sqrt{d})$ factor in front of
$P^{-1} V$ is essential to its derivation: without softmax weights,
the operator no longer corresponds to a delta-rule update in the
exponential-kernel reproducing kernel Hilbert space that the paper
derives. None of these features transfer directly to the linear-time
recurrent setting, where the full $T \times T$ attention matrix is
never materialised, the complexity envelope is linear, and the value
tensor enters a recurrent state update rather than an explicit
attention sum. Chapter~\ref{ch:proposed} introduces a separate
mechanism, the Chunked Value Decorrelation, that shares the
key-similarity decorrelation principle of LUCID but is constructed
for the linear-time recurrent substrate, with chunk-local linear
systems in place of the full $T \times T$ preconditioner. The
relationship between the two is one of shared logical motivation
rather than direct adaptation.

\section{Formal expressivity bounds}\label{sec:bounds}

This section fixes the formal scope of the empirical matrix by
locating the three proposed mechanisms within the diagonal
recurrence class, rather than re-proving the known bounds from
the related-work chapter.

The position of attention-based and recurrent linear-time
architectures in the parallel-circuit hierarchy has been established
through a connected sequence of upper-bound results that the
related-work chapter discussed in detail.
Hahn~\cite{hahn2020limitations} placed the first formal limit on
hard-attention transformers, demonstrating that PARITY and Dyck-2
are not recognisable in the limit by the family. Merrill and
Sabharwal~\cite{merrill2023parallelism} extended the bound to
log-precision attention, placing it strictly inside uniform
$\mathsf{TC}^{0}$. The survey of Strobl et
al.~\cite{strobl2024formal} unifies subsequent circuit-complexity
and logic-based characterisations of transformer expressivity. The
same $\mathsf{TC}^{0}$ ceiling applies to state-space models and
linear recurrences via the result of Merrill, Petty, and
Sabharwal~\cite{merrill2024illusion}, which extends the
characterisation to the SSM family including Mamba and demonstrates
that the apparent state of these architectures does not enable
genuine state-tracking on non-Abelian permutation tasks at any
width.

The diagonal SSM class characterised formally
by~\citet{sarrof2024expressive} is the function-class home of the
three causal backbones formalised in Section~\ref{sec:backbones}.
The single-layer definition takes the form
\begin{equation}\label{eq:diagonal-ssm}
h_{t} = A(x_{t}) \odot h_{t-1} + B(x_{t}), \qquad z_{t} = \phi(h_{t}, x_{t}),
\end{equation}
with $\odot$ the elementwise product, $A(x_{t}) \in \mathbb{R}^{d}$
(or $\mathbb{C}^{d}$) the diagonal gate at step $t$, $B(x_{t})$ the
input-write, and $\phi$ a channel-mixing transformation. The source
explicitly admits the complex-diagonal extension without loss of
generality, which places the block-complex transition of
Section~\ref{sec:primitives:dho} inside the same formal class: each
$2 \times 2$ rotation-and-damping block diagonalises over
$\mathbb{C}$ to a pair of conjugate eigenvalues $r e^{\pm i \theta}$
with $r = e^{-\lambda}$, and the block-diagonal product of $K/2$
such blocks fits the elementwise form of
equation~(\ref{eq:diagonal-ssm}) under a complex
re-parameterisation.

The principal route by which a recurrent linear-time architecture
moves out of the diagonal class is the diagonal-plus-low-rank lift
introduced by DeltaNet~\cite{schlag2021deltanet} and parallelised
by Yang et al.~\cite{yang2024deltanet}, in which a rank-one
Householder update $(I - \beta_{t} k_{t} k_{t}^{\top})$ enlarges
the reachable transition group beyond what the diagonal spectrum
alone can simulate. The DeltaProduct extension
of~\citet{siems2025deltaproduct} generalises the rank-one
Householder step to a product of $n$ Householder factors per token,
which lifts the reachable transition group further within the
diagonal-plus-low-rank class and yields a tunable trade-off between
state-tracking expressivity and per-step compute. The same
diagonal-plus-low-rank parameterisation appears in
RWKV-7~\cite{peng2025rwkv7} through its time-mix block.
Architectures in this class can in principle recognise non-Abelian
permutation languages that the diagonal class cannot, and the
impossibility result of~\citet{merrill2024illusion} applies
specifically to the diagonal class without bounding the
diagonal-plus-low-rank class above.

The three mechanisms proposed in Chapter~\ref{ch:proposed} operate
within the diagonal-class function family across all three causal
backbones. The Multi-Scale Depthwise Convolution mechanism of
Section~\ref{sec:primitives:msdc} operates on the input side without
modifying the transition operator. The Damped Harmonic Oscillator
mechanism of Section~\ref{sec:primitives:dho} extends the
real-diagonal transition to its block-complex closure but, by the
argument above, remains inside the diagonal class characterised
by~\citet{sarrof2024expressive}. The Chunked Value Decorrelation
mechanism of Section~\ref{sec:primitives:cvd} operates on the value
tensor without modifying the recurrence's transition either. None
of the three crosses the diagonal-to-low-rank boundary.

The formal results recalled here function as scope conditions
rather than theorems the thesis attempts to violate. The matrix
does not predict non-Abelian state-tracking gains at any scale,
and its probes (LibriSpeech, Common Voice, MQAR) are chosen
deliberately so as not to require them. State-tracking probes
such as $S_{n}$ permutation composition and Dyck-$k$ for
$k \ge 2$ are deferred to future work on diagonal-plus-low-rank
mechanisms; the natural extension of the framework is to repeat
the cross-architecture transfer matrix on a
diagonal-plus-low-rank backbone class with mechanisms that
target the third axis introduced in Section~\ref{sec:axes}.
